\chapter{Trade and household extensions}\label{chap:trade}

\section{Home consumption}\label{sec:home}

Households in many African economies consume part of what they produce. A SAM
records this as a payment from the household's column to the activity's row,
$\SAM(a, h)$, at producer prices. The model treats it as a separate good:
\begin{itemize}
\item the activity's output $\V{QA}_a$ is marketed output, which yields
commodities through $\theta_{ac}$ \eqref{eq:Q6}, plus home consumption
$\sum_h \V{QHA}_{ah}$;
\item home consumption is valued at the activity price $\V{PA}_a$, so it bears
no sales tax, VAT, import competition or margins;
\item it enters household demand \eqref{eq:I5H} with its own marginal budget
share $\beta^h_{ah}$ and, with LES demand, subsistence quantity
$\gamma^h_{ah}$; consumption spending $\V{EH}_h$ \eqref{eq:I4} includes it.
\end{itemize}

\textbf{Calibration.} $\V{QHA}^0_{ah} = \SAM(a, h)$, with $\V{PA}^0_a = 1$;
marketed output is $\V{QA}^0_a - \sum_h \SAM(a, h)$, which must be positive,
and $\theta_{ac}$ is $\SAM(a, c)$ divided by it. $\V{EH}^0_h$ includes home
consumption, and $\beta^h_{ah} = \SAM(a, h) / \V{EH}^0_h$. With LES demand,
home consumption of the output of activity $a$ takes the income elasticity of
$a$'s paired commodity, and
$\gamma^h_{ah} = \V{QHA}^0_{ah} + \beta^h_{ah}\,\V{EH}^0_h / (\V{PA}^0_a\,\phi_h)$.
The CPI keeps its weights on marketed consumption; welfare measures use the
full expenditure function, with $\V{PA}_a$ for home goods. GDP includes home
consumption at $\V{PA}_a$ (real GDP at $\V{PA}^0_a$). Population growth in
dynamic runs scales $\gamma^h$ as it scales $\gamma$.

\section{Export demand}\label{sec:exportdemand}

By default the country is small in world markets: it sells any quantity at the
world price $\V{pwe}_c$. For commodities given an elasticity of export demand
$\eta^e_c > 0$ (\manref{cge_calibrate}, option \cmd{exportdemand()}), the world
price is a variable $\V{PWE}_c$, and exports fall as it rises
\eqref{eq:T10}. The CET (\eqref{eq:T1}, \eqref{eq:T2}) still determines how
much producers offer at the export price $\V{PE}_c$ \eqref{eq:P2}. A larger
$\eta^e_c$ approaches the small-country case, and a change in $\V{pwe}_c$
shifts world demand. Export revenue \eqref{eq:S3}, export taxes
\eqref{eq:I8} and nominal GDP use $\V{PWE}_c$; real exports are valued at base
world prices.

\section{Export-only commodities}\label{sec:exportonly}

When a commodity's domestic output is all exported, the SAM shows exports
equal to domestic output (net of export taxes and export margins), and no
domestic sales from domestic output. By default such a commodity falls into
$CR$, since its exports are at least its domestic output
(section~\ref{sec:cal-trade}): its output goes to the domestic market and its
exports are drawn from composite supply at $\V{PQ}_c$. With
\cmd{cge_calibrate, exportonly()} its output is exported directly instead:
$\V{QD}_c$ and $\V{PD}_c$ do not exist, the CET is replaced by \eqref{eq:T3E},
and the export price $\V{PE}_c$ is the producer price. If the commodity is also
imported, its composite supply is imports alone \eqref{eq:T7}; otherwise it
has no composite supply, and $\V{PQ}_c$ and $\V{QQ}_c$ do not exist.
Calibration lists the commodities whose exports equal domestic output, so that
the choice is explicit. An export-only commodity may pay export taxes.

\section{Margins by market}\label{sec:margins-market}

Margins may differ between domestic sales, imports and exports, as in the
IFPRI standard model \citep{lofgren2002}. A margin account can be declared as
paid per unit of domestic sales (\manref{cge_map}, option
\cmd{domesticmargins()}), of imports (\cmd{importmargins()}) or of exports
(\cmd{exportmargins()}); accounts in \cmd{marginaccounts()} keep the
treatment per unit of composite supply. With $\V{MG}_{tc} = \SAM(t, c)$ and
the shares in which account $t$ buys margin services,
$s_{c't} = \SAM(c', t) / \sum_{c''} \SAM(c'', t)$:
\begin{itemize}
\item $\V{icd}_{c'c} = \sum_{t \in T^D} s_{c't}\,\V{MG}_{tc} / \V{QD}^0_c$:
domestic buyers pay $\V{PDD}_c$ (\eqref{eq:P3}, \eqref{eq:T5}), and producers
receive $\V{PD}_c$;
\item $\V{ici}_{c'c} = \sum_{t \in T^M} s_{c't}\,\V{MG}_{tc} / \V{QM}^0_c$:
the margins are part of the import price $\V{PM}_c$ \eqref{eq:P1};
\item $\V{ice}_{c'c} = \sum_{t \in T^E} s_{c't}\,\V{MG}_{tc} / \V{QE}^0_c$:
exporters receive the world price, less export taxes, less margins
\eqref{eq:P2}.
\end{itemize}
Export margins are part of the export value in the SAM, so
\begin{equation*}
\V{QE}^0_c = \SAM(c, row) - \SAM(tax_{exp}, c) - \sum_{t \in T^E} \V{MG}_{tc},
\qquad \V{pwe}_c = \SAM(c, row) / \V{QE}^0_c,
\end{equation*}
which gives $\V{PE}^0_c = 1$. The base prices
$\V{PDD}^0_c = 1 + \sum_{c'} \V{icd}_{c'c}$ and
$\V{PM}^0_c = 1 + \V{tm}_c + \sum_{c'} \V{ici}_{c'c}$ enter the Armington shares
(section~\ref{sec:cal-cet}) in place of $\V{PD}^0_c$ and $\V{PM}^0_c$.
Composite supply $\V{QQ}^0_c$ excludes export margins. A commodity that pays
domestic, import or export margins must have domestic sales, imports or
exports respectively.
